\documentclass[../../../include/open-logic-section]{subfiles}

\begin{document}

\olfileid[fr]{sth}{replacement}{finiteaxiomatizability}
\olsection{Annexe : axiomatisabilité finie}

Terminons ce chapitre en tirant quelques résultats du remplacement.
Le premier est dû à \citet{Montague1961} ; remarquons qu'il ne s'agit
pas d'une démonstration \emph{dans} $\ZF$, mais d'une démonstration
\emph{à propos de} $\ZF$ :
\begin{thm}\ollabel{zfnotfinitely}
	Si $\ZF$ est cohérente, elle n'est pas finiment axiomatisable. Plus généralement,
	si $\Th{T}$ est fini et si $\Th{T} \Proves \ZF$,
	alors $\Th{T}$ est incohérent.

	(Ici, nous nous limitons tacitement aux énoncés du premier ordre
	dont le seul symbole primitif non logique est $\in$, et nous
	écrivons $\Th{T} \Proves \ZF$ pour signifier que
	$\Th{T} \Proves \phi$ pour tout $\phi \in \ZF$.)
\end{thm}
\begin{proof}
	Fixons un $\Th{T}$ fini tel que $\Th{T} \Proves \ZF$.
	Ainsi, $\Th{T}$ démontre la réflexion, c'est-à-dire le
	\olref[sth][replacement][ref]{reflectionschema}.
	Puisque $\Th{T}$ est fini, nous pouvons en faire une
	conjonction unique, $\theta$. En appliquant la réflexion
	à cette formule,
	$\Th{T} \Proves \exists \beta(\theta \liff
	\theta^{V_\beta})$. Comme, trivialement,
	$\Th{T} \Proves \theta$, nous obtenons
	$\Th{T} \Proves \exists \beta\ \theta^{V_\beta}$.

	Soit maintenant $\psi(X)$ l'abréviation de :
	\[
		\theta^X \land X\text{ est transitif} \land
		(\forall Y \in X)(Y\text{ est transitif}
		\lif \lnot \theta^{Y})
	\]
	En gros, cela signifie que $X$ est un modèle transitif
	de $\theta$, minimal pour $\in$ parmi ces modèles.
	Puisque $\Th{T} \Proves \exists \beta\
	\theta^{V_\beta}$, les propriétés élémentaires des
	rangs dans $\ZF$, donc dans $\Th{T}$, donnent :
	\begin{equation}
		\Th{T} \Proves \exists M \psi(M).
		\tag{*}\label{Mpsi}
	\end{equation}
	Par le premier conjonctif de $\psi(X)$, chaque fois que
	$\Th{T} \Proves \sigma$, on a
	$\Th{T} \Proves \forall X(\psi(X) \lif \sigma^X)$.
	Par conséquent, d'après \eqref{Mpsi} :
	\begin{align*}
		\Th{T} &\Proves \forall X(\psi(X)
		\lif (\exists N \psi(N))^X)\\
		\intertext{À partir de là et, de nouveau, de \eqref{Mpsi} :}
		\Th{T} &\Proves \exists M(\psi(M)
		\land (\exists N \psi(N))^M)
		\intertext{Il s'ensuit en particulier que :}
		\Th{T} &\Proves \exists M(\psi(M)
		\land (\exists N \in M)((N\text{ est transitif})^M
		\land (\theta^N)^M))
		\intertext{Puis, par un raisonnement élémentaire sur la transitivité :}
		\Th{T} &\Proves \exists M(\psi(M)
		\land (\exists N \in M)(N\text{ est transitif}
		\land \theta^N))
	\end{align*}
	$\Th{T}$ est donc incohérent.\footnote{Ce
	«~raisonnement élémentaire~» exige de démontrer
	certains faits d'«~absoluité~» pour les ensembles
	transitifs. Note de l'édition française : la relativisation extérieure de la transitivité de $N$ est à $M$, et non à $N$ comme dans l'original ; cet indice est corrigé.}
\end{proof}

Voici un résultat analogue, signalé par
\citet[223]{Potter2004} :

\begin{prop}\ollabel{finiteextensionofZ}
	Soit $\Th{T}$ une extension de $\Z$ obtenue par
	l'ajout d'un nombre fini de nouveaux axiomes.
	Si $\Th{T} \Proves \ZF$, alors $\Th{T}$ est
	incohérent. (Nous adoptons ici les mêmes restrictions
	tacites que dans le \olref{zfnotfinitely}.)
\end{prop}
\begin{proof}
	Notons $\theta$ la conjonction de tous les axiomes
	de $\Th{T}$, à l'exception des instances (en
	nombre infini) de la séparation. En définissant
	$\psi$ à partir de $\theta$ comme dans le
	\olref{zfnotfinitely}, mais en remplaçant «~transitif~» par «~transitif et potent~» dans ses deux occurrences, nous pouvons montrer que
	$\Th{T} \Proves \exists M \psi(M)$.

	Comme dans le \olref{zfnotfinitely}, nous pouvons
	établir le schéma suivant : chaque fois que
	$\Th{T} \Proves \sigma$, on a
	$\Th{T} \Proves \forall X(\psi(X) \lif \sigma^X)$.
	L'absoluité de la potence de $N \in M$, lorsque $M$ est transitif et potent, permet alors de conclure
	comme dans le \olref{zfnotfinitely}.

	Cependant, établir ce schéma demande un peu plus
	de travail que dans le \olref{zfnotfinitely}. En effet,
	les instances de la séparation appartiennent à
	$\Th{T}$, mais ne sont pas des conjonctifs de
	$\theta$. Nous pouvons surmonter cet obstacle
	en démontrant, pour chaque instance $\sigma$
	de la séparation, que
	$\Th{T} \Proves \forall X(X\text{ est transitif et potent}
	\lif \sigma^X)$. Nous laissons cette vérification
	au lecteur.\footnote{Note de l'édition française : la transitivité seule, exigée dans l'original, ne suffit pas à hériter de la séparation. Par exemple, l'ordinal $3$ est transitif et contient $2$, mais pas le sous-ensemble $\{1\}$ de $2$. La preuve et l'exercice sont donc restreints aux ensembles transitifs et potents, et $\psi$ est ajustée en conséquence.}
\end{proof}
\begin{prob}
	Montrez que, pour chaque instance $\sigma$ de la
	séparation, on a
	\[
	  \Z \Proves \forall X(X\text{ est transitif et potent}
	  \lif \sigma^X)\,.
	\]
	(Nous avons employé ce schéma
	dans la
	\olref[sth][replacement][finiteaxiomatizability]{finiteextensionofZ}.)
\end{prob}
\begin{prob}
	Montrez que, pour tout $\phi \in \Z$, on a
	$\ZF \Proves \phi^{V_{\omega+\omega}}$.
\end{prob}
\begin{prob}
	Vérifiez les autres résultats schématiques invoqués
	dans les démonstrations du
	\olref[sth][replacement][finiteaxiomatizability]{zfnotfinitely}
	et de la
	\olref[sth][replacement][finiteaxiomatizability]{finiteextensionofZ}.
\end{prob}

Comme nous l'avons remarqué dans la
\olref[sth][replacement][absinf]{sec}, cela montre
que le remplacement est strictement plus fort que le
\olref[ordinals][ordtype]{thmOrdinalRepresentation}.
Ou, plus précisément : si $\Z$ augmenté de l'énoncé
«~tout bon ordre est isomorphe à un unique ordinal~»
est cohérent, alors cette théorie ne démontre pas
au moins une instance du remplacement.

	% Par hypothèse, $\psi^{V_\alpha}$ est de la forme :
	% \[
	%	(\forall A \in V_\alpha)(\exists S \in V_\alpha)
	%	(\forall x \in V_\alpha)(x \in S \liff
	%	(\phi^{V_\alpha}(x) \land x \in A))
	% \]
	% Pour le montrer, fixons $A \in V_\alpha$. Par
	% séparation, obtenons :
	% \[
	%	S = \Setabs{x \in A}{\phi^{V_\alpha}(x)}
	% \]
	% Alors $S \in V_\alpha$, puisque
	% $S \subseteq A \in V_\alpha$, et il est clair que
	% $(\forall x \in V_\alpha)(x \in S \liff
	% (\phi^{V_\alpha}(x) \land x \in A)$.

\end{document}
